% Domain-Driven Design & CQRS section

% 1 – What is DDD?
\QuestionSlide[\CategoryBadge[DDDColor!20]{DDD \& CQRS}]<1>{What is Domain-Driven Design (DDD)?}

\begin{frame}[fragile]
  \frametitle{\AnswerTitle[\CategoryBadge[DDDColor!20]{DDD \& CQRS}]{What is Domain-Driven Design?}}

  {\footnotesize
  \textbf{DDD} is a software design approach introduced by Eric Evans that centres the model on the \textbf{business domain}. Core ideas: \textbf{Ubiquitous Language} (shared vocabulary between devs and domain experts), \textbf{Bounded Contexts} (clear model boundaries), and a rich \textbf{domain model} where behaviour lives alongside state rather than in anemic service classes.
  }
\end{frame}

% 2 – Bounded Context
\QuestionSlide[\CategoryBadge[DDDColor!20]{DDD \& CQRS}]<2>{What is a Bounded Context in DDD?}

\begin{frame}[fragile]
  \frametitle{\AnswerTitle[\CategoryBadge[DDDColor!20]{DDD \& CQRS}]{What is a Bounded Context?}}

  {\footnotesize
  A \textbf{Bounded Context} is a linguistic and logical boundary within which a domain model is consistently defined. The same word (e.g.\ \emph{Product}) can mean different things in different contexts (Catalogue vs.\ Warehouse vs.\ Billing). Each BC maps to one or more microservices and has its own data store. \textbf{Context Maps} document how BCs integrate.
  }
\end{frame}

% 3 – Aggregate
\QuestionSlide[\CategoryBadge[DDDColor!20]{DDD \& CQRS}]<2>{What is an Aggregate in DDD?}

\begin{frame}[fragile]
  \frametitle{\AnswerTitle[\CategoryBadge[DDDColor!20]{DDD \& CQRS}]{What is an Aggregate?}}

  {\footnotesize
  An \textbf{Aggregate} is a cluster of domain objects treated as a single unit for data changes. It has one \textbf{Aggregate Root} — the only entry point for external code. Invariants are enforced within the aggregate boundary; references to other aggregates use IDs only. Persist and load an aggregate as one transaction.
  }

  \begin{minted}[fontsize=\scriptsize]{csharp}
public class Order : AggregateRoot  // root entity
{
    private readonly List<OrderLine> _lines = [];
    public IReadOnlyList<OrderLine> Lines => _lines.AsReadOnly();

    public void AddLine(ProductId productId, int qty, Money price)
    {
        if (qty <= 0) throw new DomainException("Qty must be positive");
        _lines.Add(new OrderLine(productId, qty, price));
        AddDomainEvent(new OrderLineAdded(Id, productId, qty));
    }
}
  \end{minted}
\end{frame}

% 4 – Value Object
\QuestionSlide[\CategoryBadge[DDDColor!20]{DDD \& CQRS}]<1>{What is a Value Object in DDD?}

\begin{frame}[fragile]
  \frametitle{\AnswerTitle[\CategoryBadge[DDDColor!20]{DDD \& CQRS}]{What is a Value Object?}}

  {\footnotesize
  A \textbf{Value Object} has no identity — two VOs with the same attributes are equal. They are \textbf{immutable} (replacing, not mutating). Use VOs for domain concepts like \texttt{Money}, \texttt{Address}, \texttt{Email}, \texttt{DateRange}. C\# \texttt{record} types are ideal: structural equality is generated automatically.
  }

  \begin{minted}[fontsize=\scriptsize]{csharp}
public record Money(decimal Amount, string Currency)
{
    public Money Add(Money other)
    {
        if (Currency != other.Currency)
            throw new DomainException("Currency mismatch");
        return this with { Amount = Amount + other.Amount };
    }
}

var price = new Money(9.99m, "USD");
var tax   = new Money(0.80m, "USD");
var total = price.Add(tax);  // new Money(10.79m, "USD")
  \end{minted}
\end{frame}

% 5 – Domain Events
\QuestionSlide[\CategoryBadge[DDDColor!20]{DDD \& CQRS}]<2>{What are Domain Events?}

\begin{frame}[fragile]
  \frametitle{\AnswerTitle[\CategoryBadge[DDDColor!20]{DDD \& CQRS}]{What are Domain Events?}}

  {\footnotesize
  \textbf{Domain Events} represent something significant that happened in the domain: \texttt{OrderPlaced}, \texttt{PaymentFailed}. Aggregates raise events; after the transaction is saved, handlers react (send email, update read model, publish integration event). This keeps business rules inside the aggregate while decoupling side effects.
  }

  \begin{minted}[fontsize=\scriptsize]{csharp}
public record OrderPlaced(OrderId OrderId, CustomerId CustomerId,
                           IReadOnlyList<OrderLine> Lines,
                           DateTimeOffset PlacedAt) : IDomainEvent;

// Handler (MediatR INotificationHandler)
public class SendOrderConfirmation(IEmailService email)
    : INotificationHandler<OrderPlaced>
{
    public Task Handle(OrderPlaced e, CancellationToken ct)
        => email.SendConfirmationAsync(e.CustomerId, e.OrderId);
}
  \end{minted}
\end{frame}

% 6 – CQRS
\QuestionSlide[\CategoryBadge[DDDColor!20]{DDD \& CQRS}]<2>{What is CQRS?}

\begin{frame}[fragile]
  \frametitle{\AnswerTitle[\CategoryBadge[DDDColor!20]{DDD \& CQRS}]{What is CQRS?}}

  {\footnotesize
  \textbf{CQRS} (Command Query Responsibility Segregation) separates \textbf{Commands} (write, change state) from \textbf{Queries} (read, never change state). This enables: independent scaling of reads/writes, optimised read models (flat DTOs, cached views), and a clean separation of domain logic (commands) from reporting logic (queries).
  }

  \begin{minted}[fontsize=\scriptsize]{csharp}
// Command
public record PlaceOrderCommand(CustomerId CustomerId,
                                 List<CartLine> Lines) : IRequest<OrderId>;

// Query
public record GetOrderSummaryQuery(OrderId OrderId)
    : IRequest<OrderSummaryDto>;

// MediatR dispatches to the correct handler
OrderId id = await _mediator.Send(placeCmd);
OrderSummaryDto dto = await _mediator.Send(new GetOrderSummaryQuery(id));
  \end{minted}
\end{frame}

% 7 – MediatR
\QuestionSlide[\CategoryBadge[DDDColor!20]{DDD \& CQRS}]<2>{What is MediatR and how does it support CQRS in .NET?}

\begin{frame}[fragile]
  \frametitle{\AnswerTitle[\CategoryBadge[DDDColor!20]{DDD \& CQRS}]{What is MediatR?}}

  {\footnotesize
  \textbf{MediatR} is a lightweight in-process mediator library. Requests (\texttt{IRequest<T>}) are dispatched to exactly one handler; notifications (\texttt{INotification}) fan out to many handlers. \textbf{Pipeline behaviours} wrap every request: logging, validation (FluentValidation), transactions, and caching can be added without touching handlers.
  }

  \begin{minted}[fontsize=\scriptsize]{csharp}
// Pipeline behaviour example: automatic validation
public class ValidationBehaviour<TReq, TRes>(IEnumerable<IValidator<TReq>> validators)
    : IPipelineBehavior<TReq, TRes> where TReq : IRequest<TRes>
{
    public async Task<TRes> Handle(TReq req, RequestHandlerDelegate<TRes> next, CancellationToken ct)
    {
        var failures = validators.SelectMany(v => v.Validate(req).Errors).ToList();
        if (failures.Count > 0) throw new ValidationException(failures);
        return await next();
    }
}
  \end{minted}
\end{frame}

% 8 – Event Sourcing
\QuestionSlide[\CategoryBadge[DDDColor!20]{DDD \& CQRS}]<3>{What is Event Sourcing?}

\begin{frame}[fragile]
  \frametitle{\AnswerTitle[\CategoryBadge[DDDColor!20]{DDD \& CQRS}]{What is Event Sourcing?}}

  {\footnotesize
  \textbf{Event Sourcing} stores the aggregate's complete history as an immutable sequence of events instead of its current state. Current state is derived by replaying events. Benefits: full audit trail, temporal queries, ability to rebuild read models. Typically paired with CQRS. Implementations: EventStoreDB, Marten (PostgreSQL), Azure Cosmos DB Change Feed.
  }

  \begin{minted}[fontsize=\scriptsize]{csharp}
// Rebuild state by replaying events
public class Order
{
    public static Order Rehydrate(IEnumerable<IDomainEvent> events)
    {
        var order = new Order();
        foreach (var e in events) order.Apply(e);
        return order;
    }

    private void Apply(IDomainEvent e) => e switch
    {
        OrderCreated c  => Status = OrderStatus.Pending,
        OrderShipped s  => Status = OrderStatus.Shipped,
        OrderCancelled  => Status = OrderStatus.Cancelled,
        _ => throw new InvalidOperationException()
    };
}
  \end{minted}
\end{frame}

% 9 – Saga / Process Manager
\QuestionSlide[\CategoryBadge[DDDColor!20]{DDD \& CQRS}]<3>{What is the Saga pattern in distributed systems?}

\begin{frame}[fragile]
  \frametitle{\AnswerTitle[\CategoryBadge[DDDColor!20]{DDD \& CQRS}]{What is the Saga pattern?}}

  {\footnotesize
  A \textbf{Saga} manages long-running distributed transactions without a two-phase commit. Each step publishes an event; subsequent services react. If a step fails, \textbf{compensating transactions} undo previous steps. Two styles:
  \begin{itemize}
    \item \textbf{Choreography} – each service decides what to do on receipt of an event (no central coordinator).
    \item \textbf{Orchestration} – a Saga Orchestrator (Process Manager) explicitly drives the sequence (e.g.\ MassTransit Sagas, NServiceBus Sagas).
  \end{itemize}
  }
\end{frame}
